\section{Variables Have Roles: Fixed, Arbitrary, Chosen, and Bound}
\chapterline

In Chapter~2 (\emph{Mathematical Sentences: Statements, Predicates, and Claims}), we saw that a predicate is an open sentence containing empty slots, like $P(x) \colon \text{``}x\text{ is prime''}$. To evaluate whether the sentence is true or false, those slots must be filled by concrete objects or bound by quantifiers.

In elementary algebra, a variable often represents an unknown number in an equation. You solve $3x + 4 = 19$, isolate $x$, and find $x = 5$, where $x$ represents a specific numerical value.

In proof-based mathematics, variables serve several distinct purposes. Letters in a proof represent different concepts, carry distinct logical permissions, and operate within specific scopes. Conflating these roles is a frequent source of error in proofs.

Consider a beginner trying to prove that every integer $n \ge 2$ has a prime factor:
\begin{quote}
\itshape Flawed attempt: Let $n = 12$. Since $12 = 3 \times 4$ and $3$ is prime, $n$ has a prime factor. Therefore, every integer $n \ge 2$ has a prime factor.
\end{quote}
The arithmetic is correct ($12$ really is $3 \times 4$), yet the argument is invalid. Why? Because when a claim begins ``for every integer $n$,'' the argument must hold for an \emph{arbitrary} integer. One cannot substitute a specific chosen number to prove a universal statement.

To read and write proofs with confidence, you need to recognize which of the four fundamental roles a variable plays in any given sentence.

\begin{historybox}
For most of human history, even elementary algebra was purely rhetorical (written out entirely in words without mathematical symbols). Medieval mathematicians had no notation for unknowns or coefficients; every problem was phrased verbally using terms like \emph{cosa} (the Italian word for ``the thing'') and \emph{census} (``the square''). In 1591, the French mathematician Fran\c{c}ois Vi\`ete made a major leap in \emph{In artem analyticem isagoge} by systematically using letters: vowels ($A, E, I$) for unknown quantities and consonants ($B, C, D$) for known parameters. For the first time, mathematicians could manipulate entire families of equations at once.

Yet introducing letters created a stubborn misconception that lasted for two centuries. In the 1700s, Leonhard Euler, Jean le Rond d'Alembert, and Daniel Bernoulli often treated variables as physical quantities that ``flow with time''---an intuition inherited from Isaac Newton's calculus of fluxions (Newton's term for the rates of change of quantities flowing continuously through time). This confusion erupted in the famous vibrating string controversy of 1747--1755, where d'Alembert, Euler, and Bernoulli argued bitterly over the motion of a plucked violin string. Because they lacked a precise definition of a function and had no formal rules for whether a letter was an arbitrary fixed input or an actively moving quantity, they frequently treated a variable as fixed in one step and varied it in the next, falling into contradictory results.

Between 1889 and 1895, the Italian mathematician Giuseppe Peano, together with Gottlob Frege in 1879, dismantled this confusion by formalizing variable binding and lexical scope (the bounded region of text where a variable name is defined and valid). In the 1930s, Alonzo Church introduced the $\lambda$-calculus (lambda calculus, a formal mathematical system where variables are explicitly bound as inputs to functions), which established the scoping rules used by modern programming languages. A variable is neither an unknown number nor a moving physical particle; it is a scoped placeholder governed by strict visibility permissions. Distinguishing the different roles of variables keeps proofs logically sound and avoids variable capture (where an expression accidentally binds a variable with the same name from an outer scope).

Today, an active research challenge in computer science and mathematical logic is how computer proof assistants (such as Lean, Coq, and Isabelle) should manage bound variables when verifying massive libraries of mathematics. Naive substitution causes severe performance bottlenecks. Designing representations that handle variable binding without variable capture---such as De Bruijn indices (numbering placeholders by how many nested scopes enclose them rather than naming them with letters) or nominal logic---remains a lively subject.
\end{historybox}

\subsection{The four roles of variables}

\begin{center}
\begin{tikzpicture}[
    font=\small,
    lane/.style={draw=#1!70!black, fill=#1!5, rounded corners=3pt, thick, text width=12.5cm, inner sep=7pt},
    badge/.style={font=\scriptsize\bfseries\scshape, text=white, fill=#1!80!black, rounded corners=2pt, inner sep=3pt}
]

% Lane 1: Fixed Constant
\node[lane=gray] (l1) at (0, 7.5) {
    \textbf{\color{gray!90!black}1. Fixed Constant} \quad \colorbox{gray!75}{\color{white}\scriptsize\bfseries Scope: Global to Problem} \hspace{0.4em} \colorbox{gray!75}{\color{white}\scriptsize\bfseries Permission: Read-Only}\\[0.3em]
    \footnotesize
    \textbf{Initiation:} ``Let $c > 0$ be given'' or ``Fix an integer $n \ge 2$.''\\
    \textbf{Behavior:} Handed down by the problem statement. Its value never changes during the argument, but you cannot assume it equals any specific number like $3$ or $10$.
};

% Lane 2: Arbitrary Universal
\node[lane=sourceblue] (l2) at (0, 5.0) {
    \textbf{\color{sourceblue!90!black}2. Arbitrary Universal} \quad \colorbox{sourceblue!80}{\color{white}\scriptsize\bfseries Scope: Universal Block} \hspace{0.4em} \colorbox{sourceblue!80}{\color{white}\scriptsize\bfseries Permission: Handed by Adversary}\\[0.3em]
    \footnotesize
    \textbf{Initiation:} ``Let $x \in \R$ be arbitrary.''\\
    \textbf{Behavior:} Introduced to prove a $\forall x$ statement. Think of $x$ as chosen by a skeptical adversary trying to break your claim. You may assume \emph{only} that $x$ belongs to the stated domain.
};

% Lane 3: Chosen Witness
\node[lane=targetgreen] (l3) at (0, 2.5) {
    \textbf{\color{targetgreen!90!black}3. Chosen Witness} \quad \colorbox{targetgreen!80}{\color{white}\scriptsize\bfseries Scope: Local Witness} \hspace{0.4em} \colorbox{targetgreen!80}{\color{white}\scriptsize\bfseries Permission: Actively Built by Prover}\\[0.3em]
    \footnotesize
    \textbf{Initiation:} ``Choose $y = x + 1$'' or ``Set $k = 2m$.''\\
    \textbf{Behavior:} Introduced to prove an $\exists y$ claim. You hold complete freedom to engineer this object, and it may depend on any variables introduced earlier in the scope.
};

% Lane 4: Bound Dummy
\node[lane=structurepurple] (l4) at (0, 0.0) {
    \textbf{\color{structurepurple!90!black}4. Bound / Dummy Variable} \quad \colorbox{structurepurple!80}{\color{white}\scriptsize\bfseries Scope: Internal Expression} \hspace{0.4em} \colorbox{structurepurple!80}{\color{white}\scriptsize\bfseries Permission: Invisible Outside}\\[0.3em]
    \footnotesize
    \textbf{Initiation:} Appears inside $\sum_{i=1}^n$, $\int_a^b f(t)\,dt$, $\{x \in \R : x > 0\}$, or $\forall x P(x)$.\\
    \textbf{Behavior:} A purely internal placeholder. It cannot be referenced outside its containing expression. Renaming it changes nothing: $\sum_{i=1}^n i = \sum_{k=1}^n k$.
};

\end{tikzpicture}\\[0.4em]
{\small\textbf{The Four Roles of Variables:} Every letter in a proof serves one of four operational roles. Its category defines where it is visible (scope), who chooses it (permissions), and whether its name can be freely swapped without altering meaning.}
\end{center}

\subsection{Lexical scope: mathematical proofs as programs}

If you have written software, you are already comfortable with variable scope (the region of source code where a name is recognized). A local variable declared inside a function body ceases to exist once that function returns. Trying to access that variable outside the curly braces causes a compile-time error.

Mathematical proofs obey the exact same scoping rules:

\begin{center}
\begin{tikzpicture}[font=\small]

% Code block comparison
\node[draw=gray!40, fill=gray!3, rounded corners=3pt, text width=5.1cm, inner sep=7pt, align=left] (code) at (-3.7, 0) {
    \textbf{\footnotesize In Software (C / Java / Rust):}\\[0.3em]
    \texttt{\footnotesize void verifyForAll(Domain D) \{}\\
    \quad\texttt{\footnotesize for (Object x : D) \{}\\
    \quad\quad\texttt{\footnotesize Object y = construct(x);}\\
    \quad\quad\texttt{\footnotesize assert(satisfies(x, y));}\\
    \quad\texttt{\footnotesize \}}\\
    \quad\texttt{\footnotesize // Error: 'y' and 'x' are out of scope!}\\
    \texttt{\footnotesize \}}
};

% Math block comparison
\node[draw=sourceblue!40, fill=sourceblue!3, rounded corners=3pt, text width=5.1cm, inner sep=7pt, align=left] (math) at (3.7, 0) {
    \textbf{\footnotesize In a Mathematical Proof:}\\[0.3em]
    \footnotesize \textbf{Goal:} Prove $\forall x \in D, \exists y \in D, P(x, y)$.\\[0.3em]
    \quad \textbf{Step 1:} Let $x \in D$ be arbitrary. \emph{(Scope opens)}\\
    \quad \textbf{Step 2:} Choose $y = f(x)$.\\
    \quad \textbf{Step 3:} Verify $P(x, y)$ holds.\\
    \quad \emph{(Scope closes)}\\
    \footnotesize In the next paragraph, $x$ and $y$ have no meaning!
};

\draw[<->, {Latex[length=2.5mm]}-{Latex[length=2.5mm]}, thick, quotientorange] (code.east) -- node[above=3pt, font=\scriptsize\bfseries, text=quotientorange, align=center]{Exact Scope\\Equivalence} (math.west);

\end{tikzpicture}\\[0.4em]
{\small\textbf{Lexical Scoping in Software and Proof:} In both programming and pure mathematics, local variables declared inside a block vanish as soon as the block terminates. Attempting to use $x$ or $y$ after the universal claim is proved is an out-of-scope error.}
\end{center}

When a proof says ``Let $x \in \R$ be arbitrary,'' it opens a local scope block. Inside that block, $x$ is treated as an immutable object about which nothing is known except that $x \in \R$. Once you show that the desired property holds for that generic $x$, the universal statement $\forall x \in \R, \ldots$ is established, and the scope block closes. You cannot subsequently refer to ``that $x$'' in later arguments.

This scoping discipline protects against a frequent proof error: the leaked variable. Suppose you are asked to prove a property about two even integers, $n$ and $m$. By definition, $n$ is even if there exists an integer $k$ such that $n = 2k$. A common beginner mistake is to write:
\begin{quote}
``Since $n$ is even, $n = 2k$. Since $m$ is even, $m = 2k$.''
\end{quote}
Here, the letter $k$ was a local witness tied specifically to the evenness of $n$. Reusing the exact same variable name for $m$ accidentally forces $n = m$, meaning you are no longer proving the claim for two arbitrary even integers, but only for the special case where both integers are equal! To maintain proper variable scoping, every fresh existential claim requires a fresh variable name: $n = 2k$ and $m = 2j$ for integers $k$ and $j$.

\begin{recoverybox}
\textbf{Mental check: Spotting the scope leak}
A student writes the following argument:
\begin{quote}
``To show that $A \subseteq B$, let $x \in A$ be arbitrary. Since $A \subseteq B$, it follows that $x \in B$. Hence $A \subseteq B$ is proved. Therefore, we have found an element $x$ that belongs to both $A$ and $B$, which proves that $A \cap B \neq \emptyset$.''
\end{quote}
What is the fatal flaw in the student's reasoning?
\\[0.3em]
\emph{Answer:} The variable $x$ was introduced as an arbitrary element strictly inside the temporary scope used to prove $A \subseteq B$. Once $A \subseteq B$ is established, that scope closes and $x$ vanishes. The argument does not show that $A$ contains any elements at all! If $A = \emptyset$, the subset claim $\emptyset \subseteq B$ is vacuously true, yet $\emptyset \cap B = \emptyset$. Claiming that $x \in A \cap B$ in the next sentence attempts to use an out-of-scope variable that has already expired.
\end{recoverybox}

\subsection{Free versus bound variables}

Look at the following three mathematical expressions:
\[
\text{(A)}\quad \sum_{i=1}^n i^2, \qquad\qquad \text{(B)}\quad \int_0^1 x^2 \, dx, \qquad\qquad \text{(C)}\quad x^2 + y^2 = 1.
\]

In expression (A), the index $i$ is \emph{bound} by the summation. The letter $n$, however, is \emph{free}: the numerical value of the sum depends directly on what $n$ is. If you change $n$ from $3$ to $4$, the sum changes from $14$ to $30$. But if you replace every occurrence of $i$ with $j$, getting $\sum_{j=1}^n j^2$, the value is completely unchanged. The index $i$ is merely an internal dummy counter.

In expression (B), the variable $x$ is bound by the integral. The expression evaluates to the concrete number $1/3$. There is no $x$ remaining in the answer. Writing ``let $x = 10$'' before the integral has zero effect on its outcome.

In expression (C), both $x$ and $y$ are free. Without additional context, the equation does not assert a truth value; it defines a predicate (an open condition satisfied by pairs of points on the unit circle).

The distinction between free and bound variables is straightforward: \textbf{free variables are parameters determined by the surrounding context, while bound variables are internal placeholders local to an operator} (Figure~\ref{03:fig-bound-variable-capsule}).

\begin{figure}[htbp]
\centering
\begin{tikzpicture}[
    font=\small,
    capsule/.style={draw=structurepurple!70, fill=structurepurple!4, rounded corners=6pt, line width=1.1pt},
    freebox/.style={draw=sourceblue!70, fill=sourceblue!8, rounded corners=4pt, thick, align=left, inner sep=6pt},
    boundbox/.style={draw=structurepurple!80, fill=structurepurple!12, rounded corners=4pt, thick, align=center, inner sep=5pt},
    outbox/.style={draw=targetgreen!70, fill=targetgreen!8, rounded corners=4pt, thick, align=left, inner sep=6pt}
]

% Left: Free inputs box
\node[freebox, font=\scriptsize, text width=3.0cm] (freenote) at (-5.3, 0) {
    \textbf{\color{sourceblue!90!black}Free Variables: $n, x$}\\[0.2em]
    \textbf{External Dials:}\\
    $\bullet$~$n$: sets loop limit\\
    $\bullet$~$x$: scales each term\\[0.3em]
    Changing either dial changes the final value.
};

% Capsule container
\draw[capsule] (-2.5, -2.1) rectangle (2.5, 1.9);
\node[anchor=north, font=\scriptsize\bfseries\scshape, text=structurepurple!85!black] at (0, 1.8) {The Bound Variable Capsule};

% Expression inside capsule
\node[font=\large] (expr) at (0, 0.75) {$\displaystyle \sum_{k=1}^n (x \cdot k)$};

% Bound note inside capsule
\node[boundbox, font=\scriptsize, text width=4.3cm] (bnote) at (0, -1.05) {
    \textbf{\color{structurepurple!90!black}Bound Variable: $k$ (Dummy Index)}\\[0.2em]
    Trapped strictly inside the capsule.\\
    Renaming $k \to j$ leaves value unchanged.\\
    Cannot be referenced outside!
};

% Arrow into capsule
\draw[-{Latex[length=2.5mm]}, thick, sourceblue!80] (freenote.east) -- node[above=2pt, font=\tiny\bfseries, text=sourceblue!90!black]{Inputs} (-2.5, 0);

% Right: Output box
\node[outbox, font=\scriptsize, text width=3.0cm] (outnote) at (5.3, 0) {
    \textbf{\color{targetgreen!90!black}Evaluated Value:}\\[0.2em]
    $\displaystyle \frac{n(n+1)}{2} \cdot x$\\[0.4em]
    \textbf{Depends only on $n, x$.}\\
    The dummy index $k$ has completely vanished!
};

% Arrow out of capsule
\draw[-{Latex[length=2.5mm]}, thick, targetgreen!80] (2.5, 0) -- node[above=2pt, font=\tiny\bfseries, text=targetgreen!90!black]{Yields} (outnote.west);

\end{tikzpicture}
\caption{Free and bound variables in an expression. Free variables ($n, x$) are parameters supplied from the surrounding context. The bound variable ($k$) is an internal index local to the summation; it cannot be referenced in the outer scope and disappears upon evaluation.}
\label{03:fig-bound-variable-capsule}
\end{figure}

Because a bound variable is purely internal, you are always free to rename it to any other symbol that is not already in use. In mathematical logic, this principle is called $\alpha$-equivalence (alpha-equivalence, meaning two formulas are identical up to renaming of bound variables):
\[
\int_a^b f(t) \, dt = \int_a^b f(u) \, du = \int_a^b f(w) \, dw.
\]
The choice of letter $t, u, w$ carries zero mathematical information. 

However, you must never rename a bound variable to a letter that is already free in the surrounding context. In the sum $\sum_{k=1}^n k$, renaming $k$ to $n$ produces $\sum_{n=1}^n n$. This name collision triggers \emph{variable capture}: the upper limit $n$ is captured by the summation index, destroying the original formula.

\begin{recoverybox}
\textbf{Mental check: Free or bound?}
In each of the following expressions, identify which letters are free and which are bound:
\begin{enumerate}[label=(\alph*)]
\item $\displaystyle \int_1^x \frac{1}{t} \, dt = \ln(x)$.
\item $\displaystyle \{y \in \R : y^2 < c\}$.
\item $\displaystyle \lim_{h \to 0} \frac{f(x+h) - f(x)}{h}$.
\end{enumerate}
\emph{Answers:} In (a), $x$ is free (it is the upper limit of integration); $t$ is bound by $dt$. In (b), $c$ is free (a parameter); $y$ is bound by the set-builder braces. In (c), $x$ and $f$ are free (the point and the function); $h$ is bound by the limit operator.
\end{recoverybox}

\subsection{A worked example of quantifier and role parsing}

Consider the statement:
\[
\forall x \in \R, \, \exists y \in \R \text{ such that } x + y = 0.
\]

Let us trace the roles and permissions line by line:
\begin{enumerate}
\item \textbf{The universal player moves first}: The statement begins with $\forall x \in \R$. Therefore, in our proof, we must begin by accepting an \emph{arbitrary} real number:
\begin{quote}
``Let $x \in \R$ be arbitrary.''
\end{quote}
We cannot choose $x = 0$ or $x = 5$. We must leave $x$ completely general.
\item \textbf{The existential player responds}: The next quantifier is $\exists y \in \R$. We must now produce a \emph{chosen witness}. Because $y$ comes \emph{after} $x$ in the quantifier order, our choice of $y$ is legally permitted to depend on $x$:
\begin{quote}
``Choose $y = -x$.''
\end{quote}
\item \textbf{Verify the requirements}: We verify two things: first, that $-x \in \R$ (it is, since the real numbers are closed under negation), and second, that $x + y = x + (-x) = 0$.
\item \textbf{Conclude and close scope}: Both conditions are satisfied, so the statement is proved, and the local variables $x$ and $y$ close their scopes.
\end{enumerate}

Now observe what happens if the quantifiers are reversed:
\[
\exists y \in \R \text{ such that } \forall x \in \R, \, x + y = 0.
\]
Here, $y$ must be chosen \emph{first}, independently of $x$. Any proposed value of $y$ (say $y = 0$) fails for $x = 1$, since $1 + 0 = 1 \neq 0$. In fact, for any fixed real number $y$, choosing $x = -y + 1$ yields $x + y = (-y + 1) + y = 1 \neq 0$. Reversing the quantifier order changes the dependence between the variables.

\begin{trapbox}
Never define an existential witness using a variable that has not yet entered scope. In the statement $\forall x \, \exists y$, your formula for $y$ may use $x$. In $\exists y \, \forall x$, your formula for $y$ must be an absolute constant chosen before $x$ is announced. Writing ``Choose $y = -x$'' before $\forall x$ has appeared is an illegal forward reference.
\end{trapbox}

\subsection{The bridge to quantifiers: who moves first?}

Scoping rules tell us where a variable is visible and what operations are legally permitted on it. But in advanced mathematics, theorems rarely involve a single variable in isolation. They involve chains of variables interacting with one another.

When you prove a theorem in analysis or algebra, you often handle several variables at once: a fixed space, an arbitrary element, a chosen witness, and a bound index. Who chooses which letter, and which choices may depend on earlier ones?

In Chapter~4 (\emph{Quantifiers: The Game of For Every and There Exists}), we examine how quantifiers ($\forall$ and $\exists$) govern these choices as a game, and why their order is critical to mathematical statements.

\subsection{Exercises}

\begin{enumerate}
\item \textbf{Role Identification (Warm-up):} In the statement
\[
\forall n \in \Z, \, \exists m \in \Z \text{ such that } n = 2m + 1,
\]
identify the roles of $n$ and $m$. Is the statement true or false? Notice that even if a statement is mathematically false, its variables still possess well-defined syntactic roles!

\item \textbf{Multi-Variable Parsing (Calculus):} In the definition of continuity of a function $f \colon \R \to \R$ at a point $c$:
\begin{quote}
``For every $\eps > 0$, there exists $\delta > 0$ such that for all $x \in \R$, if $|x - c| < \delta$, then $|f(x) - f(c)| < \eps$,''
\end{quote}
identify the exact role of each letter: $c$, $\eps$, $\delta$, and $x$. Which variable is a fixed constant, which are arbitrary universals, and which is a chosen witness? On which prior variables is the formula for $\delta$ permitted to depend?

\item \textbf{Bound Dummies and Evaluation:} In the expression
\[
L = \lim_{n \to \infty} \sum_{k=1}^n \frac{k}{n^2}:
\]
\begin{enumerate}[label=(\alph*)]
\item Which variables are bound dummies, and what binding operator captures each of them?
\item Compute the value of the expression by first summing $\sum_{k=1}^n k = \frac{n(n+1)}{2}$ and then evaluating the limit as $n \to \infty$. Does any variable survive into the final number?
\end{enumerate}

\item \textbf{Diagnosing Proof Bugs (Illegal Specialization):} Explain the exact error in the following pseudo-proof:
\begin{quote}
\itshape Claim: Every integer is even.\\
False proof: Let $n \in \Z$ be arbitrary. Choose $n = 4$. Since $4 = 2(2)$, $n$ is even. Therefore every integer is even.
\end{quote}
Which role permission was violated? Why does setting $n = 4$ destroy the universal validity of the proof?

\item \textbf{Diagnosing Proof Bugs (Forward Referencing):} Explain what is wrong with the following proof attempt:
\begin{quote}
\itshape Claim: $\exists y \in \R \text{ such that } \forall x \in \R, \, x + y = x$.\\
False proof: Let $x \in \R$ be arbitrary. Choose $y = x - x = 0$. Since $x + 0 = x$, the number $y$ works for all $x$. Therefore such a $y$ exists.
\end{quote}
Although the claim is actually true (the number $0$ indeed works for all $x$), why is this proof invalid as written? (Look at the quantifier order and the scope-opening line).

\item \textbf{Variable Capture in Calculus:} Consider the function $F(x) = \int_0^x (x - t) \, dt$.
\begin{enumerate}[label=(\alph*)]
\item Evaluate $F(x)$ by integrating with respect to $t$.
\item A student reasons: ``Because $t$ is a dummy variable, I can rename it to $x$. Then the integrand becomes $(x - x) = 0$, so $\int_0^x 0 \, dx = 0$.'' Explain why this reasoning is completely false, identifying the exact collision between a free variable and a bound variable.
\end{enumerate}

\item \textbf{Proof Lab (Writing Scopes Explicitly):} Write a rigorous, step-by-step proof of the claim:
\[
\forall x \in \R, \, (x > 3 \implies x^2 - 9 > 0).
\]
For each line of your proof, explicitly state whether you are opening a scope, introducing an arbitrary element, assuming a hypothesis, or closing a scope.
\end{enumerate}

\begin{walkbox}
When reading an advanced textbook, pause at each variable. Ask yourself: ``Is this a fixed parameter, an arbitrary element of a set, a chosen witness satisfying an existential claim, or a bound index local to an expression?''
\end{walkbox}

\newpage
